\section{Continuation and the ordinary constraints}
\label{cont:section}

The fixed-component comparison has converted the hypothesis on $B$
into equations on one spectral branch of the ancestor theory of $W$.
We now continue those equations to every branch. The essential
finiteness comes from ancestors: their cotangent powers are bounded
by the dimension of a moduli space of curves, independently of the
fiber degree. We then add the branches belonging to $X$ and use the
same fixed-component comparison in reverse.

\subsection{Why the equations are identities of germs}

Write the connected ancestor potential at background zero as
\[
 \log\mathcal A_W=\sum_{g\geq0}\hbar^{g-1}\mathcal F_g.
\]
All coefficients below are taken in the polynomial equivariant
basis of Section~\ref{geo:section}. In particular we have not
localized the invariants themselves in order to define them.

\begin{lemma}[Polynomial coefficients and ancestor bounds]
\label{cont:polynomiality}
Fix a genus $g$, an actual base class $\beta$, and a list of
insertions $b_i\bar\psi_i^{a_i}$, $1\leq i\leq n$.
Their connected invariant in $W$, summed over fiber degree with
weight $y^l$, is polynomial in $y$ and $\lambda$.
It vanishes unless the distinguished curve is stable and
\begin{equation}
                   \sum_i a_i\leq 3g-3+n.
 \label{cont:ancestor-bound}
\end{equation}
The latter bound is independent of $\beta,l,\lambda$.
\end{lemma}

\begin{proof}
By definition, curve-unstable ancestor terms are omitted. For stable
data, the product of cotangent classes is pulled back from
$\overline{\mathcal M}_{g,n}$, whose complex dimension is $3g-3+n$.
This proves \eqref{cont:ancestor-bound}, before integration and
without using equivariant degree.

For a class $(\beta,l)$, the virtual dimension of the stable-map
space is
\[
 (1-g)(\dim_\C W-3)+n+\int_\beta\kappa+(r+1)l,
 \qquad \kappa=c_1(TB)+c_1(E).
\]
The insertion degrees are fixed. Proper equivariant pushforward
takes values in $\C[\lambda]$; it has no negative cohomological
degree. Consequently the integral vanishes when this virtual
dimension exceeds the total insertion degree. Since $l\geq0$,
only finitely many $l$ can occur. Each of the remaining integrals
is polynomial in $\lambda$, which proves the claim.
This dimension argument uses total cohomological degree; the
first Hodge grading used to define the modes is a separate
grading.
\end{proof}

The modes on the $B$ branch are infinitesimal symplectic
formally at $y=0$ by Proposition~\ref{fixed:equivalence}.
By Proposition~\ref{spec:zero}, their entries are finite
polynomials in $\log y$ with coefficients holomorphic at
zero. Taylor expansion is injective on this ring, so the
symplectic identities hold as germs on a punctured sector.
This property
continues entrywise along any path on which the modes continue:
it is a linear identity between their matrix entries and their
adjoints. For a branch and logarithm obtained by such
continuation, put
\begin{equation}
 \mathcal E_{k,i}
   =\mathcal A_W^{-1}\op(A_{k,i})\mathcal A_W .
 \label{cont:equation}
\end{equation}
This notation means that the differential operator acts on the
potential, followed by multiplication by its inverse.
For the entire $X$ cluster, the notation $\mathcal E_{k,X}$
means the same expression with $A_{k,X}=\sum_{i\in X}A_{k,i}$,
using a single logarithm throughout the cluster. Projective
satisfaction means that every coefficient of positive insertion
degree in \eqref{cont:equation} vanishes.

\begin{lemma}[Finite coefficient tests]
\label{cont:finite-tests}
Fix $k\geq-1$, a power $\hbar^{G-1}$, a base class $\beta$, and an
insertion monomial of degree $N$ in \eqref{cont:equation}.
Its coefficient is a finite sum of holomorphic germs on the
spectral covering of the generic $(y,\lambda)$ locus, with the
chosen logarithms. It continues along every path in that locus.
Near $y=0$, the coefficient for $\mathcal E_{k,B}$ is a finite
polynomial in $\log y$ with coefficients holomorphic through
zero; the coefficient for the cluster expression
$\mathcal E_{k,X}$ is holomorphic through zero. In these two cases its
Taylor expansion is the corresponding formal Novikov coefficient
test.
\end{lemma}

\begin{proof}
Let $F=\log\mathcal A_W$ and use positive loop coordinates $q$.
A normally ordered quadratic operator has the schematic form
\[
 \frac{1}{2\hbar}C(q,q)
       +\sum_{\alpha,\gamma}B_{\alpha\gamma}q^\gamma
                    \partial_\alpha
       +\frac{\hbar}{2}
          \sum_{\alpha,\gamma}D_{\alpha\gamma}
                    \partial_\alpha\partial_\gamma .
\]
The tensors have the graded symmetry appropriate to their indices.
Dividing its action on $e^F$ by $e^F$ replaces its derivatives by
$\partial_\alpha F$ and by the graded expression
\[
 \partial_\alpha\partial_\gamma F+
                  (\partial_\alpha F)(\partial_\gamma F).
\]
The dilaton translation $q=t-z1$ adds only a fixed constant to
one coordinate and does not affect finiteness.

At $\hbar^{G-1}$, the first-derivative term involves
$\mathcal F_G$. The second-derivative term involves
$\mathcal F_{G-1}$, and the product of first derivatives involves
the finitely many pairs $\mathcal F_g,\mathcal F_h$ with
$g+h=G$. Terms with negative genus are absent. The multiplication
term occurs only at $\hbar^{-1}$. For a fixed output monomial,
the connected potentials which can occur have at most $N+2$
distinguished marks. Lemma~\ref{cont:polynomiality} therefore
bounds every differentiated descendant index uniformly,
independently of the fiber degree.

There is also a finite bound on the relevant matrix entries of
$A_{k,i}$. By Proposition~\ref{spec:modes} it has loop powers
at least $-1$ and finite order in $z\partial_z$.
In the second-derivative part, both negative-mode indices have
just been bounded by the ancestor inequality. In the mixed
part, the differentiated index is bounded and the multiplied
index is either an index of the prescribed output monomial or
the dilaton index. In the multiplication part, passage from
a nonnegative mode to a negative one, with lower loop bound
$-1$, permits only finitely many indices. Thus each of the
three parts uses only finitely many loop coefficients of
the operator. The $z$ derivatives multiply these entries by
polynomials in their mode indices and do not change this
conclusion.

At fixed $\beta$, Novikov finiteness leaves only finitely many
decompositions into the base classes carried by the operator
and the one or two connected potentials. The latter
coefficients are polynomials in $y$ by
Lemma~\ref{cont:polynomiality}. The relevant operator entries
are holomorphic germs that continue with the spectral branches,
by Proposition~\ref{spec:modes}. The pairing and its inverse
are fixed rational matrices in $\lambda$; we work at a nonzero
$\lambda$ germ. This proves the finite-sum assertion and
continuation.

Finally, Proposition~\ref{spec:zero} gives holomorphic
coefficients at zero for the $X$ cluster and a polynomial
dependence on $\log y$ for each $B$-branch mode. All operations
in the coefficient test have just been shown to be finite.
Taylor expansion therefore commutes with that test and yields
exactly the formal identity used in
Proposition~\ref{fixed:equivalence}.
\end{proof}

The last statement is what allows formal identities to start
analytic continuation. A finite sum
$\sum_{j=0}^J f_j(y)(\log y)^j$, with $f_j$ holomorphic at zero,
whose formal series in $y$ and $\log y$ vanishes is zero on a
punctured sector: every Taylor coefficient of every $f_j$
vanishes. Thus we use convergent germs of individual coefficient
tests, without requiring convergence of a total potential.

\subsection{Transport from one branch to the other fixed component}

Fix $\lambda\neq0$. The branch equation is
\begin{equation}
               y=h^r(h+\lambda).
 \label{cont:cover}
\end{equation}
Its critical values are $0$ and
$(-r)^r\lambda^{r+1}/(r+1)^{r+1}$. Removing these values from
the $y$-plane gives an unramified covering of degree $r+1$.

\begin{lemma}
\label{cont:monodromy}
The monodromy of \eqref{cont:cover} acts transitively on its
$r+1$ branches.
\end{lemma}

\begin{proof}
The inverse image of the complement of the critical values is
the complex $h$-plane with finitely many points removed. It is
connected and path connected. Given any two points over a
chosen regular value, join them by a path in this inverse
image. Its projection is a loop at that value, whose lift
takes the first point to the second. This is transitivity.
\end{proof}

\begin{proposition}
\label{cont:projective-transfer}
If the ordinary descendant potential of $B$ satisfies its
ordinary modes projectively, then the same is true for $X$.
\end{proposition}

\begin{proof}
By Proposition~\ref{fixed:equivalence}, the hypothesis gives all
the projective equations for $A_{k,B}$ on $\mathcal A_W$ near
$y=0$. Lemma~\ref{cont:finite-tests} turns their formal
coefficient identities into identities of germs on a
punctured sector. Their infinitesimal symplectic identities
hold there as well, by the same fixed-component comparison.
Symplecticity is an entrywise linear identity, so it
continues together with the operators.

Choose a regular value in this sector. Every coefficient of
each projective equation continues along any loop based
there. The ancestor coefficients themselves return unchanged,
because at fixed base class they are polynomials in $y$.
The only change is in the continued spectral branch and
logarithm of its mode operator. By
Lemma~\ref{cont:monodromy}, the continued $B$ branch can be
any of the $r+1$ branches. The projective equations and
infinitesimal symplecticity therefore hold on each branch.

Continuing a logarithm may add $2\pi\mathrm i$ times an
integer. Proposition~\ref{spec:modes} gives an invertible
triangular binomial change among all modes when this happens.
Consequently the simultaneous equations hold for any chosen
local logarithm on each branch.

Return to a punctured neighborhood of $y=0$ and choose the
same logarithm on the whole $X$ cluster. Its projectors are
the sum of the individual projectors, and its logarithm is
the sum of their supported logarithms. In particular
$D_X=\sum_{i\in X}D_i$. These operators contain $z$
derivatives, so it is useful to record why taking their
powers still respects this sum. Two-sided support and
commutation of each projector with $z$ give
\[
 (zD_i)(zD_j)=zD_iP_i zP_jD_j=0\qquad(i\neq j).
\]
It follows that $A_{k,X}=\sum_{i\in X}A_{k,i}$ for every
$k\geq-1$. Quantization is linear, and a sum of quantities
independent of the insertion variables is again independent
of them. Hence $\mathcal A_W$ satisfies all cluster modes
projectively.

Proposition~\ref{spec:zero} extends the cluster operators
holomorphically through $y=0$. Taking Taylor coefficients in
the equations is legitimate by
Lemma~\ref{cont:finite-tests}. We may therefore apply
Proposition~\ref{fixed:equivalence} in the reverse direction,
with $F=X$. Its conclusion concerns precisely the ordinary
descendant potential $Z_X$, with the full curve labels.
\end{proof}

\subsection{The prescribed scalar normalization}

The use of projective equations has absorbed the scalar ambiguity
of quantization. For the ordinary Virasoro operators that
ambiguity is removed by two commutators. We give the scalar
calculation because the constant in $L_0$ is part of the
assertion.

\begin{lemma}
\label{cont:commutators}
For every smooth projective complex variety $Y$, with the
operators of Section~\ref{conv:section}, one has
\begin{equation}
 [L_{-1}^Y,L_1^Y]=-2L_0^Y,\qquad
 [L_0^Y,L_k^Y]=-kL_k^Y\quad(k\geq-1,\ k\neq0).
 \label{cont:exact-commutators}
\end{equation}
\end{lemma}

\begin{proof}
The identities $[\mu_Y,R_Y]=R_Y$ and
$[\ell_{0,Y},z]=z$ give the classical commutators
\[
 [\ell_{-1,Y},\ell_{1,Y}]=2\ell_{0,Y},
 \qquad [\ell_{0,Y},\ell_{k,Y}]=k\ell_{k,Y}.
\]
Normal ordering in our Hamiltonian sign convention reverses
the commutator sign, with a scalar from contractions between
the two opposite off-diagonal blocks of the polarization;
see \cite[Section~2]{GiventalQuant}. We compute the two
possible scalars in the full super space.

For the first commutator, multiplication by $z^{-1}$
crosses from the nonnegative to the negative polarization
only at mode zero. The nonnegative output of
$\ell_{1,Y}$ on the mode $(-z)^{-1}a$ is
\[
                    (1/4-\mu_Y^2)a
\]
in mode zero. Indeed the terms involving $R_Y$ still have
negative powers at this output. Contracting the quadratic
mode-zero terms in the two orders gives the normal-order
constant
\[
 -\frac12\str(1/4-\mu_Y^2)
  =-\frac{\chi(Y)}8+\frac12\str(\mu_Y^2)
  =-2C_Y .
\]
The parity sign can also be checked on Darboux summands.
For an even coordinate,
\[
 [q^2/(2\hbar),\hbar m\partial_q^2/2]
                =-mq\partial_q-m/2.
\]
For an odd paired plane with coordinates $\theta,\psi$,
\[
 [\theta\psi/\hbar,-\hbar m\partial_\psi\partial_\theta]
          =m(1-\theta\partial_\theta-\psi\partial_\psi).
\]
The scalar is $-m/2$ on an even line and $m$
on an odd paired plane, precisely $-\str(M)/2$ for
$M=1/4-\mu_Y^2$.
Together with the classical commutator this is
$[L_{-1}^Y,L_1^Y]=-2(\op(\ell_{0,Y})+C_Y)$.

For the second commutator with $k>0$, the only crossing of
$\ell_{0,Y}$ is its $R_Y/z$ part at mode zero. The
opposite crossing of $\ell_{k,Y}$ uses a total of $k-1$
factors of $R_Y$: expanding
$z^{-1}(z\ell_{0,Y})^{k+1}$, a term that raises loop
power by one must have exactly this many zero-power
$R_Y$ factors. Its contraction with the first crossing
therefore raises first Hodge degree by $k$. Such an
endomorphism has zero trace and zero supertrace.
For $k=-1$ both crossings have the same direction, so
there is no contraction. Thus the second commutator has
no scalar correction. The constant $C_Y$ commutes with
every operator, and translating $q=t-z1_Y$ preserves
all commutators. This proves
\eqref{cont:exact-commutators}.
\end{proof}

\begin{corollary}
\label{cont:exact}
Projective satisfaction of all the ordinary modes by $Z_Y$
implies $\V(Y)$ with exactly the constant specified in
$L_0^Y$.
\end{corollary}

\begin{proof}
Projective satisfaction says $L_k^Y Z_Y=c_k Z_Y$, where
$c_k$ is independent of all insertion variables.
The operators differentiate neither the Novikov parameters
nor $\hbar$, so they commute with every $c_j$.
Their commutators consequently annihilate $Z_Y$.
The first identity in \eqref{cont:exact-commutators}
gives $L_0^Y Z_Y=0$. The second gives every remaining
equation in the required range.
\end{proof}

\begin{proof}[Proof of Theorem~\ref{thm:main}]
The hypothesis $\V(B)$ gives projective satisfaction on $B$.
By Proposition~\ref{cont:projective-transfer}, it passes to $X$.
Corollary~\ref{cont:exact} gives the prescribed ordinary
Virasoro equations.

All constructions used the full cohomology and the
categorical inverse pairings. Curve splittings, the shift
correspondence, and the final Taylor expansion preserved
the actual integral homology labels. Thus the conclusion
holds for every genus and curve class, with arbitrary
ordinary descendant insertions, as asserted.
\end{proof}
